\section{Implementation, step by step}

This is the order in which the project was actually built.

\subsection{Step 1: Review the mid-semester prototype}
The prototype (now in \texttt{legacy/}) used 5 sensors and random data. Reviewing it, we found the seed bug: the
correlation graph was estimated on layout seed 42 but the scheduler was evaluated on seed 123. Because the seed also
placed the sensors, the graph described a different set of sensors (correlation $r=0.04$ with the truth). Decision:
rebuild as a clean, tested Python package and keep the legacy code only for reference.

\subsection{Step 2: Simulator, environment and Kalman sink}
\begin{itemize}
  \item \texttt{config.py}: one \texttt{Scenario} object holds every setting ($M$, $C$, length scale, $\rho$, layout,
        noise, interference, pilot length, graph threshold, freshness bound).
  \item \texttt{simulator.py}: the correlated field (Section 5.2), with separate layout and noise seeds, and
        \texttt{fit\_model}.
  \item \texttt{graphs.py}: correlation and interference graphs.
  \item \texttt{estimator.py}: the Kalman sink.
  \item \texttt{env.py}: the slot loop, ages, the freshness guard and safety assertions, and per-slot metrics. It can
        \emph{rewind} the field so every policy sees identical readings.
\end{itemize}
Tests written at this stage: the field has the target covariance, the graph matches the sensors it is used on, the
Kalman variance is calibrated, polling reduces neighbours' uncertainty.

\subsection{Step 3: Baselines and the evaluation harness}
\texttt{policies/baselines.py} implements the five baselines and the oracle; \texttt{selection.py} has the shared
greedy and feasibility routines; \texttt{evaluate.py} runs episodes in parallel across CPU cores and summarises
metrics with confidence intervals. A parametrised test runs every baseline and checks it never breaks freshness,
capacity or interference.

\subsection{Step 4: The GNN and how its design evolved}
\begin{enumerate}
  \item First version: a GNN that scored sensors, trained by behaviour cloning on Kalman-greedy's choices. It
        under-performed its own teacher.
  \item Final version: the residual pointer decoder (Section 5.8). It starts as Kalman-greedy exactly, so learning only
        has to add improvements.
\end{enumerate}
Files: \texttt{gnn/features.py} (13 features, two graphs, batching), \texttt{gnn/model.py} (encoder, decoder, critic),
\texttt{policies/gnn\_policy.py} (wraps the model as a scheduler).

\subsection{Step 5: PPO training}
\texttt{gnn/train.py} implements Algorithm~1. Two engineering fixes mattered:
\begin{itemize}
  \item \textbf{Environment reset bug}: resetting an environment did not rewind the field, so rollouts and evaluation
        saw different readings. Fixed with \texttt{mark()}/\texttt{rewind()} and a regression test.
  \item \textbf{Long runs killed}: training stopped when the terminal session closed. We made training
        \textbf{resumable} (state saved regularly) and ran it detached; the final run was warm-started from an
        earlier checkpoint.
\end{itemize}

\subsection{Step 6: Real datasets}
Following our faculty's request for more sensors and real data, \texttt{datasets.py} adds two datasets.

\textbf{Intel Berkeley Research Lab} \cite{25}.
\begin{enumerate}
  \item Raw data: 54 Mica2Dot motes, temperature about every 31 seconds, plus mote positions in metres on a
        40\,m $\times$ 31\,m floor.
  \item Keep readings between 5 and 45\,$^\circ$C (removes sensor faults), from 28 Feb 2004 01:00 to 19 Mar 2004 (after
        that, batteries die and data degrades).
  \item Average into 5-minute slots.
  \item Keep motes with at least 70\% of slots present; fill gaps up to 2 hours (24 slots) by interpolation; drop
        motes still incomplete. Result: \textbf{45 sensors}.
  \item Chronological split: first 55\% for training (28 Feb -- 8 Mar), the rest for testing (9 -- 18 Mar).
  \item The recording has \textbf{outages} (network down, e.g.\ 10--11 and 15--16 Mar). We record where each continuous
        segment ends; test windows never cross a gap, and a live session stops cleanly at one (tested).
\end{enumerate}

\textbf{PEMS-BAY} \cite{24}.
\begin{enumerate}
  \item 325 loop detectors in the San Jose area, average speed every 5 minutes, 1 Jan -- 30 Jun 2017 (public release
        accompanying DCRNN, Zenodo record 5146275).
  \item Zero speeds mark missing data; they are interpolated.
  \item Training: January--April. Test: May--June.
\end{enumerate}

\textbf{Common preprocessing.}
\begin{itemize}
  \item \textbf{Daily profile}: for every 5-minute time of day (and for PEMS-BAY separately for weekdays and weekends),
        the average reading over the training period, smoothed over 5 neighbouring bins. The sink knows that ``rush
        hour happens''; what it must track is how \emph{today} differs from a typical day.
  \item \textbf{Standardise} the deviation from the profile by its training standard deviation.
  \item All statistics come from the training period only; a test checks that the test period never leaks into the
        model and that replayed readings match the raw data exactly.
  \item Memory: float32 storage, chunked computation, a small profile table, and pre-fitted model files reduced memory
        from about 1.75\,GB to 200--400\,MB.
\end{itemize}

\textbf{Test windows.} Intel: 10 disjoint windows of 200 slots in the test period. PEMS-BAY: 15 windows of 288 slots
(one day each). Capacity $C=6$ (Intel) and $C=40$ (PEMS-BAY), about one-eighth of the sensors.

\subsection{Step 7: Experiments and figures}
\texttt{experiments.py} runs every study and writes JSON to \texttt{results/data/}; \texttt{viz/} draws the figures in
\texttt{results/figures/}:
\begin{itemize}
  \item synthetic main benchmark (30 layouts, uniform and clustered), no-guard study, capacity / correlation / size /
        interference / pilot-length sweeps;
  \item real benchmark, real capacity sweeps, PEMS-BAY size sweep, poll-frequency map, slot-by-slot replay;
  \item significance tests and decision times.
\end{itemize}

\subsection{Step 8: Live demo}
\begin{itemize}
  \item \texttt{live.py}: a \texttt{LiveSession} replays a real test day slot by slot and runs several schedulers on
        identical readings. Each step returns the readings, each scheduler's picks (forced ones marked), the GNN's pick
        probabilities, and the sink's estimates and uncertainties.
  \item \texttt{server.py}: a small HTTP server with a JSON API: \texttt{/api/datasets}, \texttt{/api/session},
        \texttt{/api/step}, \texttt{/api/results}, \texttt{/health}.
  \item \texttt{frontend/}: plain HTML/CSS/JavaScript with Canvas: the measured field and each scheduler's view of it,
        the error over time, ``why these sensors'' bars, and a transmission log.
\end{itemize}

\subsection{Step 9: Testing, CI and release}
21 tests (\texttt{tests/test\_core.py}) run on every push with GitHub Actions. The repository is public:
\texttt{github.com/BhanuPrasadPalella-01/gnn-aoi-scheduler}.

\subsection{Code map and commands}
\begin{center}
\small
\begin{tabularx}{\textwidth}{@{}lX@{}}
\toprule
\texttt{gnn\_aoi/config.py} & Scenario settings \\
\texttt{gnn\_aoi/simulator.py} & Synthetic field, AR(1) model, \texttt{fit\_model} \\
\texttt{gnn\_aoi/datasets.py} & Intel Lab and PEMS-BAY loaders, profiles, outages, test windows \\
\texttt{gnn\_aoi/graphs.py} & Correlation and interference graphs \\
\texttt{gnn\_aoi/estimator.py} & Kalman sink \\
\texttt{gnn\_aoi/env.py} & Slot loop, AoI, freshness guard, metrics \\
\texttt{gnn\_aoi/selection.py} & Greedy variance reduction, feasible top-$C$ \\
\texttt{gnn\_aoi/policies/} & Baselines and the GNN policy \\
\texttt{gnn\_aoi/gnn/} & Features, model, PPO training \\
\texttt{gnn\_aoi/evaluate.py}, \texttt{experiments.py} & Benchmarks, sweeps, statistics \\
\texttt{gnn\_aoi/live.py}, \texttt{server.py}, \texttt{frontend/} & Live demo \\
\texttt{gnn\_aoi/viz/} & Figures and animation \\
\texttt{tests/} & 21 tests \\
\bottomrule
\end{tabularx}
\end{center}

\begin{center}
\small
\begin{tabular}{@{}ll@{}}
\toprule
\texttt{python -m gnn\_aoi download} & fetch the raw datasets \\
\texttt{python -m gnn\_aoi train} & train the GNN with PPO \\
\texttt{python -m gnn\_aoi experiments} & run all experiments \\
\texttt{python -m gnn\_aoi figures} & draw all figures \\
\texttt{python -m gnn\_aoi serve} & live demo at \texttt{http://localhost:8000} \\
\texttt{python -m pytest} & run the tests \\
\bottomrule
\end{tabular}
\end{center}
